\SourceLeaf{365}
\Needspace{12\baselineskip}
\begin{center}[2.]\end{center}

Besser auf folgende Weise:
\[
\begin{aligned}
\BareRoot{\frac{1+p+q+r}{2}}&=a\\
\BareRoot{\frac{1+p-q-r}{2}}&=b\\
\BareRoot{\frac{1-p+q-r}{2}}&=c\\
\BareRoot{\frac{1-p-q+r}{2}}&=d,
\end{aligned}
\]
\[
\begin{array}{ccc}
\tfrac12(aa+bb-cc-dd)&-ad-bc&-ac+bd\\
+ad-bc&\tfrac12(aa-bb+cc-dd)&-ab-cd\\
+ac+bd&+ab-cd&\tfrac12(aa-bb-cc+dd).
\end{array}
\]

\Needspace{13\baselineskip}
\begin{center}[3.]\end{center}

Aus der Verbindung zweier Transformationen,
\[
\left.\begin{array}{llll}
\text{deren}&\text{erster}&\text{die Scale}&a,\ b,\ c,\ d\\
\text{,,}&\text{zweiter}&\text{,,\quad,,}&\alpha,\ \beta,\ \gamma,\ \delta
\end{array}\right\}\quad\text{entspricht,}
\]
entsteht eine neue, deren Scale:
\[
\begin{gathered}
a\alpha-b\beta-c\gamma-d\delta\\
a\beta+b\alpha-c\delta+d\gamma\\
a\gamma+b\delta+c\alpha-d\beta\\
a\delta-b\gamma+c\beta+d\alpha.
\end{gathered}
\]

\Needspace{10\baselineskip}
\begin{center}[4.]\end{center}

Sind die Coordinaten des ruhenden Punkts $\xi,\eta,\zeta$, Vergrösserung $=nn$, Drehung $=\lambda$, so kann man setzen
